% !TEX program = lualatex

\documentclass[a4paper,12pt]{article}
\usepackage{fancyhdr}
\usepackage[top=0.7in,bottom=0.6in,left=0.3in,right=0.3in,headheight=15pt]{geometry}
\usepackage{graphicx}
\usepackage{hyperref}
% \usepackage{appendix}
\hypersetup{
    hidelinks
}
\tracinglostchars=1
\usepackage{setspace}
\usepackage{array}
\usepackage{gensymb}
\usepackage{pgfplots}
\usepackage{amsmath}
\usepackage{tikz}
\usepackage{circuitikz}
\usepackage{textcomp}
\usepackage{luatexja-fontspec}
\usepackage{xcolor}
\usepackage{bookmark}
\usepackage{amsmath,amssymb,amsfonts, mathrsfs}
\usepackage{tabularx}
\usepackage{booktabs}
\usepackage{longtable}
\newcolumntype{M}[1]{>{\centering\arraybackslash\(}p{#1}<{\)}}  % math-mode paragraph column for longtable (\( expands to $, but keeps VS Code highlighting sane)
\newcolumntype{L}[1]{>{\raggedright\arraybackslash}p{#1}}     % ragged-right paragraph column for longtable
\usepackage{babel}
\hbadness=10000
\definecolor{wave_dark}{RGB}{35, 110, 190}
\definecolor{wave_med}{RGB}{95, 155, 220}
\definecolor{wave_light}{RGB}{200, 225, 255}
\definecolor{raycolor}{RGB}{215, 45, 130} % Distinct magenta/pink for rays
\definecolor{dotcolor}{RGB}{0, 115, 230}  % Vivid blue for object/image points
\definecolor{rayred}{RGB}{220, 50, 50}     % Red for ray vectors
\definecolor{wfblack}{RGB}{30, 30, 30}     % Dark charcoal for incident wavefronts
\definecolor{wfblue}{RGB}{0, 100, 220}     % Royal blue for reflected wavefronts
\definecolor{wfgreen}{RGB}{30, 150, 60}    % Forest green for refracted wavefronts
\setmainjfont{FandolSong-Regular.otf}[
    BoldFont = FandolHei-Regular.otf,
    ItalicFont = FandolKai-Regular.otf,
    BoldItalicFont = FandolHei-Regular.otf
]
\setsansjfont{FandolHei-Regular.otf}
\ctikzset{bipoles/galvanometer/height/.initial=.60}
\ctikzset{bipoles/galvanometer/width/.initial=.60}

\ctikzset{bipoles/galvanometer/arrow rotate/.initial=0,% <=new key
rotation/.style={bipoles/galvanometer/arrow rotate=#1},% style for ease introduction in code
}
% Command to draw the B-field (into page) marker (an X inside a circle)
\newcommand{\BInto}[1]{
  \draw[bfieldgreen, thick] (#1) circle (0.16cm);
  \draw[bfieldgreen, thin] (#1) ++(-0.09cm, -0.09cm) -- ++(0.18cm, 0.18cm);
  \draw[bfieldgreen, thin] (#1) ++(-0.09cm, 0.09cm) -- ++(0.18cm, -0.18cm);
}
\makeatletter
\pgfcircdeclarebipole{}{\ctikzvalof{bipoles/galvanometer/height}}{galvanometer}{\ctikzvalof{bipoles/galvanometer/height}}{\ctikzvalof{bipoles/galvanometer/width}}{
    \def\pgf@circ@temp{right}
    \ifx\tikz@res@label@pos\pgf@circ@temp
        \pgf@circ@res@step=-1.2\pgf@circ@res@up
    \else
        \def\pgf@circ@temp{below}
        \ifx\tikz@res@label@pos\pgf@circ@temp
            \pgf@circ@res@step=-1.2\pgf@circ@res@up
        \else
            \pgf@circ@res@step=1.2\pgf@circ@res@up
        \fi
    \fi

    \pgfpathmoveto{\pgfpoint{\pgf@circ@res@left}{\pgf@circ@res@zero}}       
    \pgfpointorigin \pgf@circ@res@other =  \pgf@x  \advance \pgf@circ@res@other by -\pgf@circ@res@up
    \pgfpathlineto{\pgfpoint{\pgf@circ@res@other}{\pgf@circ@res@zero}}
    \pgfusepath{draw}

    \pgfsetlinewidth{\pgfkeysvalueof{/tikz/circuitikz/bipoles/thickness}\pgfstartlinewidth}

        \pgfscope
            \pgfpathcircle{\pgfpointorigin}{1.0\pgf@circ@res@up}
            \pgfusepath{draw}       
        \endpgfscope    

    \pgftransformrotate{\ctikzvalof{bipoles/galvanometer/arrow rotate}}% <= magic line
    \pgfsetlinewidth{\pgfstartlinewidth}

    \pgfsetarrowsend{latex}
    \pgfpathmoveto{\pgfpoint{0.7\pgf@circ@res@up}{\pgf@circ@res@zero}}
    \pgfpathlineto{\pgfpoint{0.7\pgf@circ@res@down}{\pgf@circ@res@zero}}
    \pgfusepath{draw}
    \pgfsetarrowsend{}
}
\def\pgf@circ@galvanometer@path#1{\pgf@circ@bipole@path{galvanometer}{#1}}
\compattikzset{galvanometer/.style = {\circuitikzbasekey, /tikz/to path=\pgf@circ@galvanometer@path}}

\ctikzset{european}
\newcolumntype{Y}{>{\raggedright\arraybackslash}X}  % ragged right inside X
\makeatletter
\renewcommand*\env@matrix[1][*\c@MaxMatrixCols c]{%
  \hskip -\arraycolsep
  \let\@ifnextchar\new@ifnextchar
  \array{#1}}
\makeatother
\pgfplotsset{width=10cm, compat=1.18}
\usetikzlibrary{datavisualization,shapes.arrows}
\usetikzlibrary{patterns,calc, positioning,arrows.meta, angles, quotes,math, 3d, shapes.misc}
\usetikzlibrary{datavisualization.formats.functions,shadings}
\usetikzlibrary{shapes.symbols,arrows.meta,decorations.pathmorphing,decorations.markings,babel,shapes.geometric,backgrounds,shadows}
\singlespacing
\pagestyle{fancy}
\fancyhf{}
\lhead{\leftmark}
% \renewcommand{\arraystretch}{2}
\renewcommand{\footrulewidth}{0.6pt}
\rfoot{\thepage}
\lfoot{\rightmark}
\newcommand{\unnumberedsection}[1]{%
  \section*{#1}%
  \phantomsection
  \addcontentsline{toc}{section}{#1}%
  \markboth{\MakeUppercase{#1}}{}%
}
\title{\textbf{IB Physics HL Notes}}
\author{Kenny Lin / BusyPizza}
